% ============================================================================
% 2.3 字符串复制：用户态的 strcpy
% ============================================================================
\subsection{copy\_string：从 input\_buf 复制到新分配的堆内存}

我们用 \texttt{sys\_read} 把键盘输入读入 \texttt{input\_buf}，
但这块缓冲区是全局共享的，下一次读入会覆盖。因此必须把当前读到的
内容复制一份到新分配的内存里。

\begin{lstlisting}[style=asm,caption=memorycopy.s — copy\_string]
.include "src/memory.inc"
.extern heap
.extern heap_end
.extern input_buf
.extern output_buf
.extern memory_alloc
.global copy_string

copy_string:
    stp x29, x30, [sp, #-32]!
    stp x19, x20, [sp, #16]

    mov x19, x0              ; x19 = 输入字节数（sys_read 的返回值）
    sub x19, x19, #1         ; 减 1，把换行符排除

    ; 先分配一块大小为 x19 的内存
    adrp x20, input_buf@PAGE
    add  x20, x20, input_buf@PAGEOFF
    bl   memory_alloc        ; 返回值 x0 是新内存块的 data 区首地址

    mov x2, #0               ; 循环计数器

copy_loop:
    cmp x2, x19              ; 是否已复制完全部字节？
    b.ge copy_done
    ldrb w4, [x20, x2]       ; 从 input_buf 读一个字节
    strb w4, [x0, x2]        ; 写到新内存
    add  x2, x2, #1
    b    copy_loop

copy_done:
    add x19, x19, #1
    add x2, x2, #1
    strb wzr, [x0, x19]      ; 在末尾写一个 0 字节（C 字符串结尾）

    ldp x19, x20, [sp, #16]
    ldp x29, x30, [sp], #32
    ret
\end{lstlisting}

\paragraph{函数输入输出}
调用者把 \texttt{sys\_read} 的返回值（实际读取的字节数）放到 \texttt{x0}，
然后 \texttt{bl copy\_string}。返回时 \texttt{x0} 指向一块新分配的内存，
里面是用户输入的内容，末尾加了一个 \\0 字节。

\paragraph{为什么要手动加 \texttt{wzr}？}
\texttt{sys\_read} 只是把字节原样从终端拷到缓冲区，不做任何处理。
它不会给你加字符串结尾标记。所以 \texttt{strb wzr, [x0, x19]} 这一步是
在告诉任何后续的"按字符串方式处理"的代码：到这里为止。

\paragraph{关于 \texttt{x19} 减 1 又加 1 的设计}
\begin{itemize}
    \item 进入时 \texttt{sys\_read} 返回的字节数包含末尾的换行
    \item \texttt{sub x19, x19, \#1} 把换行符排除在复制循环之外
    \item 但在写入终止 \texttt{0} 时，用的是原来的 \texttt{x19+1} 的位置
        （也就是换行符所在的位置），用 0 覆盖了它
\end{itemize}

这样最终得到的字符串既不包含换行，又是一个合法的 C 风格字符串。
